\documentclass[12pt]{article}
\usepackage[utf8]{inputenc}
\usepackage[T1]{fontenc}
\usepackage{hyperref}
\title{Uokha Language Resources for K-12 Curriculum Development}
\author{Compiled by Claude Sonnet 5 (claude-sonnet-5) for Andr\'as Kornai}
\date{}
\begin{document}
\maketitle

\section{Introduction}
Uokha (uok) is an Edoid entry of Edo State, Nigeria. Its status as a separate language is doubtful. The Wikipedia Edoid article says that Ihievbe and Uokha are listed in Ethnologue (2009) under a heading for spurious languages. Uokha is also the name of a town in Owan East Local Government Area, Edo State. Glottolog's bibliography for Elugbe's Edoid study mentions Uokha [uok] only in a note and gives no data for it. I could not confirm the Glottocode uokh1238 in the search results. No educational or linguistic resources for Uokha were found.

\section{Primary Educational Resources}
None found.

\section{Linguistic and Reference Materials}
\begin{itemize}
\item Wikipedia, Edoid languages (context, including the spurious-language note): \url{https://en.wikipedia.org/wiki/Edoid_languages}
\item Glottolog reference entry for Elugbe 1989 on Edoid: \url{https://glottolog.org/resource/reference/id/91279}
\end{itemize}

\section{Religious and Community Resources}
None found.

\section{Digital and Audio Resources}
None found.

\section{Curriculum Development Context}
There is no evidence of a distinct Uokha language with its own community, orthography or school provision. A first step would be to establish with Edo State communities in Owan East whether the speech is a variety of a neighbouring Edoid language such as Owan. An honest conclusion is that no curriculum resources currently exist.
\end{document}